\documentclass[12pt,a4paper]{ctexart}
\setCJKmainfont{SimSun}[BoldFont=SimHei]
\setCJKsansfont{SimHei}

\usepackage{geometry}
\geometry{left=2.5cm, right=2.5cm, top=2.5cm, bottom=2.5cm, headheight=14pt}
\usepackage{graphicx}
\usepackage{booktabs}
\usepackage{hyperref}
\usepackage{float}
\usepackage{caption}
\usepackage{enumitem}
\usepackage{fancyhdr}
\usepackage{titlesec}
\usepackage{tabularx}
\usepackage{listings}
\usepackage{xcolor}
\usepackage{setspace}
\onehalfspacing

\captionsetup[figure]{font={small,bf}, labelfont={small,bf}, textfont={small,bf}}

\pagestyle{fancy}
\fancyhf{}
\fancyhead[L]{\small 星露谷动作X — 操作说明文档}
\fancyhead[R]{\small 谢宇鹏\quad 5120243390}
\fancyfoot[C]{\thepage}
\renewcommand{\headrulewidth}{0.4pt}

\titleformat{\section}{\LARGE\bfseries\sffamily}{\thesection}{1em}{}
\titleformat{\subsection}{\Large\bfseries\sffamily}{\thesubsection}{1em}{}

\lstset{
    basicstyle=\small\ttfamily,
    numbers=left,
    numberstyle=\tiny\color{gray},
    frame=single,
    breaklines=true,
    captionpos=b,
    xleftmargin=10pt,
    xrightmargin=10pt,
}

\begin{document}

\begin{titlepage}
\thispagestyle{empty}
\centering
\vspace*{2cm}
{\Huge\bfseries\sffamily 星露谷动作X}

\vspace{0.5cm}
{\LARGE Stardew Valley Action X}

\vspace{1cm}
{\Large 操作说明文档}

\vspace{3cm}
\begin{spacing}{2.0}
\begin{tabular}{rl}
\textbf{版本：} & V11 Final \\
\textbf{开发者：} & 谢宇鹏 \\
\textbf{学号：} & 5120243390 \\
\textbf{班级：} & 软件工程2401 \\
\textbf{日期：} & 2026年6月28日 \\
\textbf{仓库：} & \url{https://github.com/xieyupeng123/shixun} \\
\end{tabular}
\end{spacing}
\vfill
\end{titlepage}

\tableofcontents
\newpage

\section{项目概述}

星露谷动作X（Stardew Valley Action X）是一款基于Python 3.12 + Pygame 2.6开发的2D俯视角动作角色扮演游戏。玩家扮演冒险者，在开放地图中探索、战斗、升级，最终击败三阶段Boss获得胜利。

\subsection{技术规格}

\begin{table}[H]
\centering
\begin{tabular}{ll}
\toprule
\textbf{项目} & \textbf{规格} \\
\midrule
开发语言 & Python 3.12 \\
游戏框架 & Pygame 2.6 \\
窗口分辨率 & 1280 × 720 像素 \\
帧率 & 60 FPS（delta time驱动） \\
地图格式 & CSV（四层叠加） \\
存档格式 & JSON（原子写入） \\
资源路径 & 全部相对路径 \\
代码总量 & 2891行（22个模块） \\
\bottomrule
\end{tabular}
\end{table}

\subsection{系统架构概览}

本项目采用四层分层架构设计，严格遵循上层依赖下层、下层绝不依赖上层的原则：

\begin{figure}[H]
\centering
\includegraphics[width=\textwidth]{../文档/代码思维导图.png}
\caption{代码整体思维导图}
\label{fig:mindmap}
\end{figure}

\begin{table}[H]
\centering
\small
\begin{tabularx}{\textwidth}{>{\raggedright\arraybackslash}p{2cm}>{\raggedright\arraybackslash}p{5cm}>{\raggedright\arraybackslash}p{5.5cm}}
\toprule
\textbf{层级} & \textbf{职责} & \textbf{核心模块} \\
\midrule
控制层 & 游戏主循环、状态机、输入事件、界面渲染 & Game, UI, Upgrade \\
逻辑层 & 关卡管理、战斗调度、Boss生命周期、地图加载 & Level, CombatManager, BossManager, MapManager \\
实体层 & 玩家、敌人、武器、魔法、粒子、瓦片 & Player, Enemy, Weapon, MagicPlayer, Tile \\
数据层 & 全局配置、资源缓存、寻路算法、存档 & settings, ResourceManager, pathfinding\_utils \\
\bottomrule
\end{tabularx}
\caption{四层分层架构}
\end{table}

\section{环境安装}

\subsection{系统要求}

\begin{itemize}
    \item 操作系统：Windows 10/11、macOS、Linux
    \item Python版本：3.10 及以上
    \item 依赖库：pygame 2.6+
\end{itemize}

\subsection{安装步骤}

\begin{enumerate}[label=\textbf{步骤\arabic*：}]
    \item \textbf{安装Python}：从 \url{https://www.python.org/downloads/} 下载并安装Python 3.12。

    \item \textbf{安装Pygame}：打开终端（命令提示符），执行以下命令：
    \begin{lstlisting}[caption={安装Pygame}]
pip install pygame
    \end{lstlisting}

    \item \textbf{获取项目}：从GitHub克隆或直接复制项目文件夹到本地：
    \begin{lstlisting}[caption={克隆项目}]
git clone https://github.com/xieyupeng123/shixun.git
    \end{lstlisting}

    \item \textbf{进入code目录}：
    \begin{lstlisting}[caption={启动游戏}]
cd stardew-valley-action-x-final/code
python main.py
    \end{lstlisting}
\end{enumerate}

\subsection{项目文件结构}

\begin{table}[H]
\centering
\small
\begin{tabular}{ll}
\toprule
\textbf{路径} & \textbf{说明} \\
\midrule
\texttt{code/main.py} & 游戏入口，全局状态机 \\
\texttt{code/settings.py} & 所有游戏配置数据 \\
\texttt{code/constants.py} & 常量定义 \\
\texttt{code/entity.py} & Entity基类 \\
\texttt{code/player.py} & Player玩家类 \\
\texttt{code/enemy.py} & Enemy敌人类 \\
\texttt{code/level.py} & Level协调者 + 三个Manager \\
\texttt{code/map\_manager.py} & 地图加载器 \\
\texttt{code/weapon.py} & 武器精灵 \\
\texttt{code/magic.py} & 魔法系统 \\
\texttt{code/particles.py} & 粒子特效系统 \\
\texttt{code/ui.py} & HUD界面渲染 \\
\texttt{code/upgrade.py} & 升级菜单 \\
\texttt{code/tile.py} & 瓦片精灵 \\
\texttt{code/resource\_manager.py} & 资源缓存单例 \\
\texttt{code/sound\_manager.py} & 音效管理器 \\
\texttt{code/music\_state.py} & BGM状态管理 \\
\texttt{code/support.py} & 工具函数 \\
\texttt{code/pathfinding\_utils.py} & A*寻路算法 \\
\texttt{code/save\_manager.py} & 存档系统 \\
\texttt{data/map/} & CSV地图数据文件 \\
\texttt{graphics/} & 图片资源 \\
\texttt{audio/} & 音频资源 \\
\texttt{font/} & 字体文件 \\
\bottomrule
\end{tabular}
\caption{项目文件结构}
\end{table}

\section{操作指南}

\subsection{游戏启动}

运行\texttt{main.py}后，首先看到主菜单界面：

\begin{itemize}
    \item 金色标题「Stardew Valley Action X」以正弦波动画浮动
    \item 蓝色粒子在背景中缓慢移动
    \item 两个菜单选项：「Start Game」和「Quit Game」
    \item 如果有存档文件，会弹出对话框询问继续或新游戏
\end{itemize}

\subsection{操控说明}

\begin{table}[H]
\centering
\small
\begin{tabularx}{\textwidth}{>{\raggedright\arraybackslash}p{3cm}>{\raggedright\arraybackslash}p{3.5cm}>{\raggedright\arraybackslash}p{6cm}}
\toprule
\textbf{操作} & \textbf{按键} & \textbf{说明} \\
\midrule
四方向移动 & W/A/S/D 或 ↑/↓/←/→ & 对角线速度自动归一化，不会加速 \\
武器攻击 & 鼠标左键 & 朝向当前方向生成武器精灵，各武器冷却不同 \\
魔法释放 & 鼠标右键 & 治疗术（回血+光环粒子）或火焰术（前方5格AOE） \\
切换武器 & Q键 或 鼠标滚轮 & 5种武器循环切换，200ms冷却防误触 \\
切换魔法 & E键 & 治疗/火焰循环切换 \\
升级菜单 & B键 & 暂停游戏，打开五维属性升级面板 \\
世界地图 & M键 & 全屏显示世界地图覆盖层 \\
存档 & P键 & 在升级菜单打开时按P存档 \\
无敌模式 & 空格键 & 激活无敌：速度×2、攻击×2、金色光效、免疫伤害 \\
\bottomrule
\end{tabularx}
\caption{全部操控按键}
\end{table}

\subsection{武器系统}

游戏提供5种武器，在\texttt{settings.py}中配置，通过Q键或鼠标滚轮循环切换：

\begin{table}[H]
\centering
\small
\begin{tabular}{lccc}
\toprule
\textbf{武器} & \textbf{基础伤害} & \textbf{冷却时间} & \textbf{特点} \\
\midrule
sword（剑） & 15 & 100ms & 均衡型，攻速快 \\
lance（枪） & 30 & 400ms & 高伤害，冷却长 \\
axe（斧） & 20 & 300ms & 中等伤害与速度 \\
rapier（细剑） & 8 & 50ms & 极速攻击，单发伤害低 \\
sai（匕首） & 10 & 80ms & 快速型，伤害略高于细剑 \\
\bottomrule
\end{tabular}
\caption{五种武器属性}
\end{table}

\textbf{伤害计算公式}：

\[
\text{总伤害} = \text{玩家攻击力}(stats['attack']) + \text{武器伤害}(weapon\_data[weapon]['damage'])
\]

无敌模式下伤害翻倍：

\[
\text{无敌总伤害} = (\text{玩家攻击力} + \text{武器伤害}) \times 2
\]

\subsection{魔法系统}

提供两种魔法，通过E键切换：

\begin{table}[H]
\centering
\small
\begin{tabular}{lccc}
\toprule
\textbf{魔法} & \textbf{强度} & \textbf{能量消耗} & \textbf{效果} \\
\midrule
heal（治疗术） & 50 & 10 & 恢复玩家生命值（不超过上限），播放治疗粒子 \\
flame（火焰术） & 10 & 20 & 在玩家前方1-5格生成火焰AOE，每格加随机偏移 \\
\bottomrule
\end{tabular}
\caption{两种魔法属性}
\end{table}

\subsection{怪物AI系统}

游戏包含4种怪物，每种具有独立的属性配置和AI行为：

\begin{table}[H]
\centering
\small
\begin{tabular}{lccccc}
\toprule
\textbf{怪物} & \textbf{生命} & \textbf{伤害} & \textbf{速度} & \textbf{探测范围} & \textbf{攻击范围} \\
\midrule
squid（鱿鱼） & 50 & 15 & 150 & 360px & 80px \\
raccoon（浣熊/Boss） & 300 & 40 & 100 & 400px & 120px \\
spirit（幽灵） & 40 & 8 & 200 & 350px & 60px \\
bamboo（竹子） & 35 & 5 & 150 & 300px & 50px \\
\bottomrule
\end{tabular}
\caption{四种怪物属性}
\end{table}

怪物AI采用五状态有限状态机，每帧根据与玩家的距离决定行为：

\begin{enumerate}[label=\textbf{状态\arabic*：}]
    \item \textbf{idle（待机）}：玩家距离 > notice\_radius，原地播放待机动画
    \item \textbf{move（追踪）}：玩家进入探测范围，使用A*算法寻路追踪
    \item \textbf{attack（攻击）}：玩家进入攻击范围，造成伤害，400ms冷却
    \item \textbf{hit（受击）}：被玩家攻击后，300ms无敌闪烁 + 击退效果
    \item \textbf{death（死亡）}：生命 ≤ 0，播放死亡粒子 + 掉落经验球
\end{enumerate}

\subsection{Boss三阶段系统}

Boss（raccoon）具有三个阶段，每击杀一次晋升下一阶段：

\begin{table}[H]
\centering
\small
\begin{tabular}{cccc}
\toprule
\textbf{阶段} & \textbf{属性倍率} & \textbf{颜色} & \textbf{实际生命值} \\
\midrule
Phase 1 & ×1.0 & 无着色（正常） & 300 \\
Phase 2 & ×2.0 & 淡红色 (255, 120, 120) & 600 \\
Phase 3 & ×4.0 & 紫色 (180, 80, 220) & 1200 \\
\bottomrule
\end{tabular}
\caption{Boss三阶段属性}
\end{table}

击杀Boss 3次后触发胜利界面。

\subsection{无敌模式}

V11版本的无敌模式特性：

\begin{itemize}
    \item \textbf{激活}：按空格键（只能激活，不能手动关闭）
    \item \textbf{持续}：30秒后自动过期
    \item \textbf{冷却}：5分钟冷却期，期间按空格无效
    \item \textbf{代价}：到期扣除当前血量50\%（至少保留1HP）+ 清零全部能量
    \item \textbf{增益}：移动速度×2，攻击伤害×2，免疫一切伤害
    \item \textbf{视觉}：角色叠加金色半透明光效
\end{itemize}

\subsection{升级系统}

按B键暂停游戏，打开升级面板。面板包含五种属性：

\begin{itemize}
    \item health（生命值）— 上限600
    \item energy（能量值）— 上限300
    \item attack（攻击力）— 上限30
    \item magic（魔法强度）— 上限15
    \item speed（移动速度）— 上限720
\end{itemize}

操作方式：←/→选择属性，↑升级（消耗经验），↓降级（退还经验）。

\subsection{存档系统}

\begin{itemize}
    \item \textbf{存档时机}：在升级面板打开时（B键暂停），按P键存档
    \item \textbf{存档内容}：玩家位置、属性、武器选择、已击杀敌人、已破坏草地
    \item \textbf{存档机制}：JSON原子写入（先写临时文件，成功后原子替换），防止写入中断导致的数据损坏
    \item \textbf{读档}：启动时自动检测\texttt{savegame.json}，存在则弹出恢复对话框
\end{itemize}

\section{游戏逻辑流程}

下图展示了游戏从启动到退出的完整执行流程：

\begin{figure}[H]
\centering
\includegraphics[width=\textwidth]{../文档/游戏逻辑流程图.png}
\caption{游戏逻辑完整流程图}
\label{fig:flow}
\end{figure}

主循环每帧执行以下步骤：
\begin{enumerate}[label=(\arabic*)]
    \item 计算delta time（帧间隔时间）
    \item 处理事件队列（键盘/鼠标/窗口事件）
    \item Player.update(dt)：输入处理 → 冷却 → 状态 → 动画 → 移动 → 能量恢复 → 无敌检测
    \item Enemy.enemy\_update(player)：所有敌人AI状态机更新
    \item player\_attack\_logic()：武器碰撞检测与伤害计算
    \item \_check\_boss\_music()：BGM区域检测（短路优化）
    \item custom\_draw()：Y轴排序渲染所有精灵
    \item ui.display()：HUD渲染
    \item 检查死亡/胜利状态
\end{enumerate}

\section{BGM音乐系统}

游戏包含三种BGM：

\begin{table}[H]
\centering
\small
\begin{tabular}{llll}
\toprule
\textbf{BGM类型} & \textbf{文件} & \textbf{音量} & \textbf{触发条件} \\
\midrule
普通BGM & game\_bgm.ogg & 0.3 & 游戏启动，默认播放 \\
Boss BGM & boss\_bgm.ogg & 0.5 & 玩家首次进入Boss的notice\_radius \\
无敌音效 & invincible.ogg & 0.5 & 按空格激活无敌模式（循环播放）\\
\bottomrule
\end{tabular}
\caption{BGM类型说明}
\end{table}

BGM状态管理采用操作系统信号量机制：MusicState类维护一个类级别状态变量（0=普通，1=Boss），所有BGM操作必须通过get()、set()、reset()三个原语进行。无敌模式结束后根据MusicState信号量值恢复正确的BGM。

\section{常见问题}

\begin{enumerate}[label=\textbf{Q\arabic*：}]
    \item \textbf{游戏启动报FileNotFoundError？}

    确保从\texttt{code/}目录运行\texttt{main.py}，所有资源路径基于此目录的相对路径。

    \item \textbf{帧率低于60FPS？}

    检查是否有其他程序占用CPU。关闭后台程序，降低粒子效果数量。

    \item \textbf{音效无声？}

    检查音频文件是否存在于\texttt{audio/}目录。音效加载失败已做降级处理，不会影响游戏运行。

    \item \textbf{存档读取失败？}

    确保\texttt{savegame.json}格式完整。可删除存档文件以新游戏开始。

    \item \textbf{如何添加新武器/怪物？}

    只需在\texttt{settings.py}的对应字典中添加配置条目，准备美术资源，无需修改核心代码。
\end{enumerate}

\section{开发者备忘}

\begin{itemize}
    \item \textbf{添加新地图}：在\texttt{data/map/}创建4个CSV文件（FloorBlocks/Grass/Objects/Entities），准备\texttt{ground.png}用于水域检测
    \item \textbf{添加新武器}：在\texttt{settings.py}的\texttt{weapon\_data}字典添加配置，在\texttt{graphics/weapons/}放置4方向PNG
    \item \textbf{添加新怪物}：在\texttt{settings.py}的\texttt{monster\_data}字典添加配置，在\texttt{constants.py}的\texttt{ENTITY\_MONSTER\_MAP}注册CSV编码
    \item \textbf{切换AI模型}：使用项目提供的\texttt{switch-model.ps1}脚本，一行命令切换模型并保留对话历史
\end{itemize}

\end{document}
